\section{Technical Architecture (`src`)}
The \texttt{src} directory has been refactored into a modular, production-ready architecture. It follows a "Source Priority" model to maximize data density across 14 seasons.

\subsection{Modular Scrapers (`src/scrapers/`)}
We transitioned from monolithic clients to a specialized package structure:
\begin{itemize}
    \item \textbf{uefa.py}: Primary source for official kickoff times and referee teams.
    \item \textbf{espn.py}: Main engine for statistical events (possession, shots, fouls) and match timelines.
    \item \textbf{flashscore.py}: High-reliability fallback for missing match statistics using internal feed APIs.
    \item \textbf{worldfootball.py}: Specialized in historical lineup recovery (2010-2015) and stadium metadata.
    \item \textbf{fbref.py}: Provides granular performance data (assists, tackles). Features a \textbf{Wayback Machine fallback} to bypass bot protections on historical records.
\end{itemize}

\subsection{Multi-Layered Enrichment Flow}
The pipeline implements an automated cascading lookup:
\begin{enumerate}
    \item \textbf{Phase 1 (API)}: UEFA and ESPN APIs are queried for structural and statistical data.
    \item \textbf{Phase 2 (Diagnostic Index)}: If data gaps persist, the system consults the pre-built indices for Flashscore and Worldfootball.
    \item \textbf{Phase 3 (Deep Fallback)}: The FBref scraper is triggered as a final attempt for high-granularity tactical metrics.
    \item \textbf{Phase 4 (Normalization)}: The \texttt{config.py} module resolves naming inconsistencies using a fuzzy-alias mapping system.
\end{enumerate}
